\documentclass[10pt,a4paper]{amsart}
\usepackage[margin=19mm]{geometry}
\usepackage{amsmath,amssymb,mathtools}
\usepackage[scr=rsfs,cal=euler]{mathalfa}
\usepackage{xfrac}
\usepackage[unicode,hidelinks,bookmarks=false]{hyperref}
\newcommand{\ie}{i.e.\ }
\newcommand{\cf}{cf.\ }
\newcommand{\resp}{respectively\ }
\newcommand{\Subsection}{\subsection}
\providecommand{\nobreakdash}{\nobreak\hbox{-}}
\setlength{\emergencystretch}{2em}
\multlinegap0pt
\usepackage{algorithm}\usepackage{algpseudocode}
\usepackage{amssymb}
\usepackage{mathtools}
\newtheorem{lemma}{Lemma}
\newtheorem{proposition}{Proposition}
\newcommand{\E}{\mathcal{E}}
\newcommand{\MaxIE}{\mathrm{MaxIE}}
\newcommand{\R}{\mathbb{R}}
\newcommand{\V}{{\mathrm{Vol}}}
\newcommand{\tr}{\mathrm{tr}}
\title{The Polarity Process for the Maximum-Volume Inscribed Ellipsoid Problem}
\author{Kaizhao Sun}
\date{Source: arXiv:2609.10888v1 (CC BY 4.0)}
\begin{document}
\maketitle

\noindent\emph{Source and licence.} The Polarity Process for the Maximum-Volume Inscribed Ellipsoid Problem. Version arXiv:2609.10888v1 (CC BY 4.0), distributed by arXiv under the Creative Commons Attribution 4.0 International licence, http://creativecommons.org/licenses/by/4.0/, as recorded in the arXiv licence metadata for this exact version and shown on the abstract page https://arxiv.org/abs/2609.10888v1. Full licence notice: \texttt{licenses/arxiv-2609.10888-CC-BY-4.0.txt}.\par\smallskip

\noindent\textit{Editorial scope.} Counted unit: the article's lemma derivative\_of\_g (source0.tex lines 886--897). The article's complete original proof of it is delivered with it (source0.tex lines 898--994, one \texttt{proof} environment of 97 lines). No mathematics is written, repaired or re-derived: the statement and the proof are the article's own lines, in the article's own notation, and no retained line is substituted or re-flowed.  Retained uncounted as the source-local context the proof needs: lines 563--575; lines 711--719; lines 724--726; lines 745--747; lines 826--837.  Nothing was omitted from any retained span, so the package carries no omission record.  No retained line contains a citation, so the wrapper carries no bibliography and no bibliography entry.  No retained line contains a figure environment, an image-inclusion command or drawing code, so no image is delivered and the package declares no asset dependency.  Editorial formatting only, in addition: an AMS \texttt{amsart} preamble that restates the article's own declarations and macros used by the retained lines, namely the environments \texttt{lemma}, \texttt{proposition} on separate counters, together with the standard packages those lines need.  The retained environments print in this document as Lemma 1, Proposition 1 in order of appearance, whereas the article prints the same items with the numbering of its own full document.  The complete native source of the article as distributed in the arXiv e-print for this exact version is bundled unchanged as \texttt{sources/209962/source0.tex}.
For $h\in\R^n$, choose $\bar t>0$ with
$x+th\in\mathrm{int}(P)$ for $t\in(-\bar t,\bar t)$, and set
\begin{align}\label{eq: x_along_h}
    x(t) := x + th,~ t\in(-\bar t,\bar t),
\end{align}
and, holding $\lambda^*\in\Lambda(x)$ fixed, parameterize
\begin{alignat}{2}\label{eq: parameterization_in_t}
     \tilde{a}_i(t) & :=  \tilde{a}_i(x(t)),\quad 
    A(t) && :=  A(x(t), \lambda^*), \notag \\ 
     c(t) &:=  c(x(t), \lambda^*),~
    S(t) && :=  S(x(t), \lambda^*).
\end{alignat}
We suppress $t$ at $t=0$; overdots denote the corresponding derivatives.

\begin{lemma}\label{lemma: differentiation}
    For $h\in\R^n$, let $x(t)$ and the associated quantities be as in
    \eqref{eq: x_along_h} and \eqref{eq: parameterization_in_t}. At $t=0$, it holds that 
    \begin{align*}
        \dot{\tilde{a}}_i = & (\tilde{a}_i^\top h)\tilde{a}_i, \quad \dot{c} = Ah, \\
        \dot{A} = & 2 \sum_{i=1}^m \lambda_i^* (\tilde{a}_i^\top h) \tilde{a}_i \tilde{a}_i^\top, \\
        \dot{S} = & 2 \sum_{i=1}^m \lambda_i^* (\tilde{a}_i^\top h) \tilde{a}_i \tilde{a}_i^\top - Ahc^\top -c h^\top A.
    \end{align*}
\end{lemma}

    \begin{align}\label{eq: differentiating_tilde_a}
        \dot{\tilde{a}}_i =\frac{a_i^\top h}{(1-a_i^\top x - ta_i^\top h)^2} a_i\Bigg|_{t=0} = \frac{a_i^\top h}{(1-a_i^\top x)^2}a_i = (\tilde{a}_i^\top h)\tilde{a}_i.
    \end{align}

\begin{proposition}\label{proposition: v_gradient}
    The function $v$ is convex and differentiable on $\mathrm{int}(P)$, with gradient given by $\nabla v(x) = (n+1)c(x)$.
\end{proposition}

\begin{lemma}\label{lemma: facts_about_g}
Let $x^*$ be the center of $\MaxIE(P)$ and let
$\lambda^*\in\Lambda(x^*)$. Then:
\begin{enumerate}
    \item $g(x) \leq v(x)$ for all $x \in \mathrm{int}(P)$, and $g(x^*) = v(x^*)$.
    \item $c(x^*)=\mathbf0$ and
    $v(x^*)=-\log\V(\MaxIE(P))$; moreover,\\
    $\MaxIE(P)=\E(x^*,Q(x^*)^{-1})$.
    \item $Q(x^*) = \frac{1}{n}S(x^*)^{-1}$.
    \item for every $i\in [m]$ with $\lambda^*_i > 0$, we have $\tilde{a}_i(x^*)^\top S(x^*)^{-1}\tilde{a}_i(x^*) = n$.
\end{enumerate}
\end{lemma}

\begin{lemma}\label{lemma: derivative_of_g}
    Given a direction $h\in \R^n$, define quantities
    \begin{align}\label{eq: g_derivative_components}
        & p_h := S(x^*)^{1/2} h, \quad
        w_i := S(x^*)^{-1/2} \tilde{a}_i(x^*), \notag \\
        & T_h := \sum_{i=1}^m \lambda^*_i (\tilde{a}_i(x^*)^\top h)(w_iw_i^\top).
    \end{align}
    The derivative of $g_h$ at $0$ satisfies $g_h'(0) = 0$, i.e., $\nabla g(x^*) = \mathbf{0}$. In addition, 
    \begin{align*}
            g''_h(0) = (3n-1)\|p_h\|^2 - 2\|T_h\|_F^2.
    \end{align*}
\end{lemma}

\begin{proof}
    Lemma \ref{lemma: differentiation} gives the first derivatives at $t=0$.
    Applying the calculation from Proposition \ref{proposition: v_gradient} to
    \begin{align}\label{eq: g_h_in_t}
        g_h(t) = \frac{1}{2}\log \det S(t) + \frac{n}{2} \log n,
    \end{align}
    its derivative at zero is
    \begin{align*}
        g'_h(0) = & \frac{1}{2}\tr(S(x^*)^{-1}\dot{S})
        = \sum_{i=1}^m \lambda_i^* (\tilde{a}_i(x^*)^\top h)\tilde{a}_i(x^*) ^\top S(x^*)^{-1}\tilde{a}_i(x^*) \\
        = & n\sum_{i=1}^m \lambda_i^* (\tilde{a}_i(x^*)^\top h) =  n \left(\sum_{i=1}^m \lambda_i^* \tilde{a}_i(x^*)\right)^\top h = n c(x^*)^\top h =0,
    \end{align*}
    where Lemma \ref{lemma: facts_about_g} gives the last line.

    For the second derivative, denote those of $(\tilde{a}_i,A,S)$ at $t=0$ by
    $(\ddot{\tilde{a}}_i,\ddot A,\ddot S)$.
    \begin{itemize}
        \item \textbf{Differentiating $\dot{\tilde{a}}_i$.} By \eqref{eq: differentiating_tilde_a}, we have $\tilde{a}'_i(t) = (a_i^\top h)(1 - a_i^\top x^*-t a_i^\top h)^{-2} a_i$, which implies
        \begin{align}\label{eq: second_order_tilde_a}
        \ddot{\tilde{a}}_i= \tilde{a}''_i(0) =  & \frac{2(a_i^\top h)^2}{(1-a_i^\top x^* - ta_i^\top h)^3} a_i \Bigg|_{t=0} \notag \\
        = & 2\left(\frac{a_i^\top h}{1-a_i^\top x^*}\right)^2 \frac{a_i}{1-a_i^\top x^*} = 2(\tilde{a}_i^\top h)^2 \tilde{a}_i.
        \end{align}
        \item \textbf{Differentiating $\dot{A}$.} Differentiating
        $A(t)=\sum_{i=1}^m\lambda_i^*\tilde a_i(t)\tilde a_i(t)^\top$ twice gives
        \begin{align*}
            A''(t) = & \sum_{i=1}^m \lambda_i^*\left(\frac{d}{dt} \tilde{a}'_i(t)\tilde{a}_i(t)^\top + \frac{d}{dt}\tilde{a}_i(t) \tilde{a}'_i(t)^\top \right) \\
            = &  \sum_{i=1}^m \lambda_i^*\left(\tilde{a}''_i(t) \tilde{a}_i(t)^\top + 2\tilde{a}'_i(t)\tilde{a}'_i(t)^\top + \tilde{a}_i(t)\tilde{a}''_i(t)^\top  \right).
        \end{align*}
        At $t=0$, Lemma \ref{lemma: differentiation} and
        \eqref{eq: second_order_tilde_a} give
        \begin{align}\label{eq: second_order_A}
            \ddot{A} = A''(0) =\sum_{i=1}^m \lambda_i^* \left(\ddot{\tilde{a}}_i \tilde{a}_i^\top + 2 \dot{\tilde{a}}_i\dot{\tilde{a}}_i^\top + \tilde{a}_i \ddot{\tilde{a}}_i^\top \right) =  6 \sum_{i=1}^m \lambda_i^* (\tilde{a}_i^\top h)^2 \tilde{a}_i \tilde{a}_i^\top.
        \end{align}
        \item \textbf{Differentiating $\dot{S}$.} Differentiating
        $S=A-cc^\top$ twice and using Lemma \ref{lemma: differentiation},
        \eqref{eq: second_order_A}, and Lemma \ref{lemma: facts_about_g} gives
        \begin{align}
            \ddot{S} = \ddot{A} - (\ddot{c}c^\top + 2\dot{c}\dot{c}^\top + c\ddot{c}^\top) =  6 \sum_{i=1}^m \lambda_i^* (\tilde{a}_i^\top h)^2 \tilde{a}_i \tilde{a}_i^\top - 2(Sh)(Sh)^\top.\label{eq: second_order_S}
        \end{align}
    \end{itemize}
    By \eqref{eq: g_h_in_t} and the fact that $g'_h(t) = \frac{1}{2}\tr(S(t)^{-1}S'(t))$, we have
    \begin{align}
     g_h''(t) = & \frac{1}{2}\frac{d}{dt}\tr(S(t)^{-1}S'(t)) = \frac{1}{2} \tr\left(\frac{d}{dt} S(t)^{-1}S'(t)\right) \notag \\
        = & \frac{1}{2} [\tr((S^{-1})'(t)S'(t)) +  \tr(S^{-1}(t)S''(t))], \label{eq: g_second_derivative}
    \end{align}
    by linearity of trace and the product rule. Differentiating $S^{-1}S=I$ gives
    \begin{alignat*}{2}
        & S(t)^{-1} S(t) = I \quad &&   \\
        \Rightarrow\quad  & (S^{-1})'S + S^{-1}S' = 0 \quad && \text{(differentiate both sides)}\\
        \Rightarrow \quad & (S^{-1})'S S^{-1}+ S^{-1}S'S^{-1}  = 0 \quad && (\text{right multiply by $S^{-1}$}) \\
        \Rightarrow \quad & (S^{-1})' = -S^{-1}S'S^{-1}.
    \end{alignat*}
    Substitution in \eqref{eq: g_second_derivative} gives
    \begin{align}
        g''_h(t)
        = \frac{1}{2}\tr(S^{-1}(t)S''(t)) - \frac{1}{2}\tr(S^{-1}(t)S'(t)S^{-1}(t)S'(t)),\notag
    \end{align}
    and hence, at $t=0$,
    \begin{align}\label{eq: g_second_derivative_2}
        g''_h(0) = \frac{1}2\tr(S(x^*)^{-1} \ddot{S}) - \frac{1}{2}\tr(S(x^*)^{-1}\dot{S}S(x^*)^{-1}\dot{S}).
    \end{align}
    We evaluate the two terms in \eqref{eq: g_second_derivative_2}, writing
    $S=S(x^*)$ and using $S^{-1}=S^{-1/2}S^{-1/2}$ and cyclicity of trace.
    \begin{itemize}
        \item By \eqref{eq: second_order_S} and part (4) of Lemma \ref{lemma: facts_about_g}, we have
        \begin{align}
            S^{-1/2}\ddot{S} S^{-1/2} =& 6 \sum_{i=1}^m \lambda^*_i (\tilde{a}_i^\top h)^2 S^{-1/2} \tilde{a}_i\tilde{a}_i^\top S^{-1/2} - 2S^{-1/2}(Sh)(Sh)^\top S^{-1/2} \notag \\
            = &  6 \sum_{i=1}^m \lambda^*_i (\tilde{a}_i^\top h)^2 w_iw_i^\top - 2p_hp_h^\top  \notag, \\
            \tr(S^{-1}\ddot{S}) = & \tr(S^{-1/2}S^{-1/2}\ddot{S} ) = \tr(S^{-1/2}\ddot{S} S^{-1/2}) \notag \\
            = & 6\sum_{i=1}^m \lambda^*_i(\tilde{a}_i^\top h)^2\|w_i\|^2 - 2\|p_h\|^2  \notag \\
            = & 6n\sum_{i=1}^m \lambda^*_i(\tilde{a}_i^\top h)^2 - 2\|p_h\|^2. \notag
        \end{align}
        Moreover, it holds that 
        \begin{align}
            &\sum_{i=1}^m \lambda^*_i(\tilde{a}_i^\top h)^2  = \sum_{i=1}^m \lambda^*_i(w_i^\top p_h)^2 = p_h^\top \left(\sum_{i=1}^m \lambda^*_i w_iw_i^\top\right) p_h,\notag \\
            &\sum_{i=1}^m \lambda^*_i w_iw_i^\top = S^{-1/2}\left(\sum_{i=1}^m \lambda_i^*\tilde{a}_i\tilde{a}_i^\top \right)S^{-1/2} = S^{-1/2}SS^{-1/2} = I_n. \notag 
        \end{align}
        As a result, we have 
        \begin{align}\label{eq: trace_S_inv_ddot_S}
            \tr(S^{-1}\ddot{S}) = 2(3n-1)\|p_h\|^2.
        \end{align}

        \item By Lemma \ref{lemma: differentiation} and the fact that
        $c=c(x^*)=\mathbf{0}$ from Lemma \ref{lemma: facts_about_g}, we have
        \begin{align*}
            S^{-1/2} \dot{S} S^{-1/2} = 2\sum_{i=1}^m \lambda_i^* (\tilde{a}_i^\top h)w_iw_i^\top = 2T_h,
        \end{align*}
        and hence 
        \begin{align}\label{eq: trace_S_inv_dot_S_sq}
            \tr(S^{-1}\dot{S}S^{-1}\dot{S})  = &   \tr(S^{-1/2}S^{-1/2}\dot{S}S^{-1/2}S^{-1/2}\dot{S}) \notag \\
            = & \tr((S^{-1/2} \dot{S} S^{-1/2} )(S^{-1/2} \dot{S} S^{-1/2})) \notag  \\
            = & \tr((2T_h)^2) = 4\tr(T_h^2) = 4\|T_h\|_F^2.
        \end{align}
    \end{itemize}
    Substituting \eqref{eq: trace_S_inv_ddot_S} and
    \eqref{eq: trace_S_inv_dot_S_sq} into \eqref{eq: g_second_derivative_2} proves the claim.
\end{proof}

\end{document}
